% !PRINTED 334
\weberpage{334}
For the necessary and sufficient condition that (2) be a basis of $\mathfrak{a}$ consists in this, that every integer $\alpha$ divisible by $\mathfrak{a}$, and only such, is contained in the form:
\weq{3}{\alpha = x_1\alpha_1 + x_2\alpha_2 + \cdots + x_n\alpha_n}
with integer $x_1, x_2, \ldots, x_n$. Consequently all numbers
\weq{4}{\beta = x_1\alpha_1\eta + x_2\alpha_2\eta + \cdots + x_n\alpha_n\eta}
are divisible by $\mathfrak{b}$. If, conversely, $\beta$ is an integer divisible by $\mathfrak{b}$, then $\beta/\eta$ is divisible by $\mathfrak{a}$ and consequently representable in the form (3). Hence $\beta$ is representable by (4).

Further, in the notation of (9) and (10), § 87, if $\alpha$ is a positive basis:
\[
\begin{aligned}
\Sigma \pm \alpha_{1,1}\alpha_{2,2}\ldots\alpha_{n,n} &= N(\mathfrak{a})\, \Sigma \pm \omega_{1,1}\omega_{2,2}\ldots\omega_{n,n}, \\
\Sigma \pm \beta_{1,1}\beta_{2,2}\ldots\beta_{n,n} &= N(\eta)\, \Sigma \pm \alpha_{1,1}\alpha_{2,2}\ldots\alpha_{n,n}, \\
&= N(\eta)N(\mathfrak{a})\, \Sigma \pm \omega_{1,1}\omega_{2,2}\ldots\omega_{n,n}, \\
&= \pm N(\mathfrak{b})\, \Sigma \pm \omega_{1,1}\omega_{2,2}\ldots\omega_{n,n},
\end{aligned}
\]
and since by (1)
\weq{5}{N(\mathfrak{b}) = N(\eta)N(\mathfrak{a})}
holds, the positive sign must stand in $\pm N(\mathfrak{b})$.

If $\lambda$ is a basis form of $\mathfrak{a}$, then $\eta\lambda$ is a basis form of $\mathfrak{b}$, and
\weq{6}{\begin{aligned} N(\lambda) &= N(\mathfrak{a})T, \\ N(\eta\lambda) &= N(\mathfrak{b})T, \end{aligned}}
and thus the same form $T$ belongs to the two ideals $\mathfrak{a}$, $\mathfrak{b}$. If we unite equivalent forms $T$ into one class, then by 2. we can say:

3. \gesperrt{To every ideal class there belongs a definite class of forms.}

In answering the question whether different ideal classes can belong to one and the same class of forms, we make the assumption that $\Omega$ is a \gesperrt{normal field}: let us assume, then, that equivalent forms $T$ and $T'$ belong to two ideals $\mathfrak{a}$ and $\mathfrak{a}'$; then, if $\lambda$ is a basis form of $\mathfrak{a}$,
\weq{7}{N(\lambda) = N(\mathfrak{a})T,}
and if we carry $T'$ over into $T$ by a linear substitution, and understand by $\lambda'$ a basis form of $\mathfrak{a}'$ (with the same variables $t$ as $\lambda$), then
\weq{8}{N(\lambda') = N(\mathfrak{a}')T.}
